\section{Identifying the mechanism}
\label{sec:identifiability}

This section shows that, under exact geometric-mean local normalization, the extrema-based
post-local block statistic used by the original policy carries no block-scale information, so
the coupling stage cannot use it, and it derives the simpler mechanism that remains.

\subsection{Collapse of post-normalization block centers}

For each nonzero local row $r$, write $m_r$ and $M_r$ for its smallest and largest
nonzero coefficient magnitudes. The geometric-mean kernel applies
$d_r=(m_rM_r)^{-1/2}$. Let $\widetilde m_r=d_rm_r$ and $\widetilde M_r=d_rM_r$ denote the
exported extrema. For a nonempty set $I$ of local rows, the block statistic is
\[
 s_I=\sqrt{\left(\min_{r\in I}\widetilde m_r\right)
                 \left(\max_{r\in I}\widetilde M_r\right)}.
\]

\begin{proposition}[Exact collapse and the effect of a skip band]
\label{prop:audit-collapse}
If every row receives its geometric-mean factor, then $s_I=1$ for every nonempty grouping
$I$ of local rows. More generally, suppose the only safeguard replaces $d_r$ by $1$ when the
candidate lies in $(1/u,u)$, where $u\geq1$. Then $s_I\in[1/u,u]$, independently of the
original coefficient scales and of the grouping of the rows.
\end{proposition}
\begin{proof}
Under exact normalization, $\widetilde m_r=\sqrt{m_r/M_r}=1/\widetilde M_r$. Consequently
$\min_r\widetilde m_r=1/\max_r\widetilde M_r$, proving the first statement.
For the second, let $c_r=\sqrt{\widetilde m_r\widetilde M_r}$. An accepted row has
$c_r=1$; a skipped row has $c_r=\sqrt{m_rM_r}\in(1/u,u)$. Thus
$u^{-2}\leq\widetilde m_r\widetilde M_r\leq u^2$ for every row. Choose $p$ attaining
the maximum exported row magnitude and $q$ attaining the minimum. Then
\[
 u^{-2}\leq\widetilde m_q\widetilde M_q
 \leq\widetilde m_q\widetilde M_p
 \leq\widetilde m_p\widetilde M_p\leq u^2.
\]
Taking square roots gives the result.
\end{proof}

For a coupling row $a_r$, let $\beta(j)$ assign column $j$ to a local block, and
define the block-weighted candidate and its ownership-blind counterpart by
\[
 \gamma_r^{\mathrm{block}}
 =\left(\sum_j a_{rj}^2/s_{\beta(j)}^2\right)^{-1/2},
 \qquad
 \gamma_r^{\mathrm{L2}}=\left(\sum_j a_{rj}^2\right)^{-1/2}.
\]
Blockless columns may use an arithmetic mean of valid block scales or the unit
fallback; both conventions preserve the interval in \cref{prop:audit-collapse}.

\begin{proposition}[Loss of ownership information]
\label{prop:audit-ownership}
Under exact local normalization and with no coefficient filtering,
$\gamma_r^{\mathrm{block}}=\gamma_r^{\mathrm{L2}}$ for every coupling row, so
changing the ownership map cannot change the candidate export; applying the same
coupling safeguards to identical candidates preserves equality.
Under the skip-band assumptions of \cref{prop:audit-collapse},
\[
 u^{-1}\leq\frac{\gamma_r^{\mathrm{block}}}
                       {\gamma_r^{\mathrm{L2}}}\leq u,
\]
and, before coupling safeguards, the condition numbers of the two exports differ by
at most a factor $u^2$ on their common nonzero singular spectrum.
\end{proposition}
\begin{proof}
The equality follows by substituting $s_i=1$. In the second case,
\[
 u^{-2}\|a_r\|_2^2\leq\sum_j a_{rj}^2/s_{\beta(j)}^2
 \leq u^2\|a_r\|_2^2,
\]
which gives the factor bound. Write the block-based export as $EA_{\mathrm{L2}}$,
where $E$ has unit entries on local rows and entries in $[1/u,u]$ on coupling
rows. Multiplication by $E$ changes each singular value by a factor between
$1/u$ and $u$ and preserves rank. Applying these inequalities to the largest and
smallest positive singular values in both directions proves the condition-number
bound.
\end{proof}

Both statements assume that the block summary sees every coefficient. Coefficient-window
vetoes and filtered extrema lie outside the hypotheses: a row with magnitudes
$(10^{-8},3,10^8)$ is geometrically centered, but restricting the summary to
$(10^{-6},10^6)$ produces $s_I=3$, a nonunit statistic created by the filter rather than by
retained physical scale. The export audit in \cref{sec:results} measures how often the
implemented safeguards break exact equality.

\subsection{What a matched comparison identifies}

With equal coverage and equal local kernels, two exports agree on their local rows, so any
difference is confined to the coupling rows. This localizes the difference without
attributing it to ownership, because the coupling kernels may still differ. The ownership
control therefore keeps local GM and coupling L2 normalization but replaces every block
scale by one (Role-Hybrid in \cref{tab:policies}); SA-Mat versus Role-Hybrid is the
ownership comparison. Flat-GM changes the coupling kernel and Flat-L2 changes the local
kernel, so neither isolates ownership.

Equality also requires equal \emph{applied} factors, not merely an inactive coupling stage in
one policy. For local rows $x_1\leq1$, $x_2\leq1$ and the coupling row
$x_1+10^{-2}x_2\leq1$, both block scales are one and the Euclidean coupling candidate
$(1+10^{-4})^{-1/2}$ falls inside the skip band $(1/2,2)$, whereas the flat GM candidate
$10$ is applied. Both policies cover every row, yet their exports differ.

\subsection{The spectral slope in a stylized model}

Take $B\geq2$ local singleton rows with positive coefficients $v_i=10^{e_i}$,
where $\min_i e_i=-R$ and $\max_i e_i=R$, and one coupling row $v^\top$.
Exact local normalization produces $I_B$ and unit post-local block centers. A
whole-row coupling factor $\gamma$ gives
\[
 A_\gamma=\begin{bmatrix}I_B\\\gamma v^\top\end{bmatrix},
 \qquad A_\gamma^\top A_\gamma=I_B+\gamma^2vv^\top,
 \qquad
 \kappa_2(A_\gamma)=\sqrt{1+\gamma^2\|v\|_2^2}.
\]
Flat geometric-mean normalization selects $\gamma=1$, whereas both the post-local block
rule and the hybrid Euclidean control select $\gamma=1/\|v\|_2$. Hence
\[
 \kappa_2(A_{\mathrm{GM}})=\sqrt{1+\sum_i10^{2e_i}},
 \qquad
 \kappa_2(A_{\mathrm{block}})=\kappa_2(A_{\mathrm{hybrid}})=\sqrt2.
\]
For fixed $B$ the ratio is of order $10^{-R}$: the advantage grows with block heterogeneity
and survives removal of the ownership map. Every coupling coefficient is multiplied by the
same $\gamma$, so all within-row ratios are preserved; dividing each $v_i$ by its own
$10^{e_i}$ would change the constraint and is not a row scaling.

\subsection{Optimal coupling factor after local whitening}

The stylized model suggests a constructive replacement for the metadata explanation: once
local scale has been removed, choose the coupling factor from the coupling coefficients
themselves. In the idealized triangular model this choice has an exact solution.

\begin{proposition}[Optimal common coupling factor after local whitening]
\label{prop:audit-triangular-optimum}
Let $R\in\mathbb R^{m\times n}$, with $m,n\geq1$, and set $\rho=\|R\|_2$. For
\[
 C_\gamma=\begin{bmatrix}I_n&0\\\gamma R&\gamma I_m\end{bmatrix},
 \qquad\gamma>0,
\]
the unique factor minimizing $\kappa_2(C_\gamma)$ is
\[
 \gamma_*=(1+\rho^2)^{-1/2},
 \qquad
 \min_{\gamma>0}\kappa_2(C_\gamma)=\sqrt{1+\rho^2}+\rho.
\]
On an admissible interval $0<\ell\leq\gamma\leq h$, the optimum is
$\min\{h,\max\{\ell,\gamma_*\}\}$.
\end{proposition}
The proof (Online Resource 1) reduces the spectrum by an SVD to two-by-two blocks. For the
block at $\rho$, the singular-value ratio $k$ satisfies
$k+k^{-1}=\gamma^{-1}+\gamma(1+\rho^2)$, whose unique minimizer gives the result.

For a single coupling row, $\gamma_*$ is the Euclidean normalization of the complete row,
including the coefficient of its aggregate variable. The idealized coupling problem
therefore has a simple optimal whole-row normalization that uses no ownership information,
and this is the mechanism the matched control isolates empirically. The statement concerns a
common factor on a locally whitened triangular structure; for general generated MIPs it
motivates the Euclidean coupling kernel rather than certifying it.

\paragraph{Why conditioning needs a residual bound.}
The interval form of \cref{prop:audit-triangular-optimum} matters. In $A_\gamma$,
$\kappa_2$ decreases to one as $\gamma\to0$, while a fixed residual on the exported coupling
row corresponds to an original residual larger by the factor $1/\gamma$. A conditioning
criterion alone cannot say how small a coupling factor may be: on any interval its
minimizer for $A_\gamma$ is the lower endpoint $\ell$. \Cref{sec:contract} derives that
endpoint from an application-unit residual budget.
